%%%%%%%%%%%%%%%%%%%%%%%%%%%%%%%%%%%%%%%%%%%%%%%%%%%%%%%%%%%%%%%%%%%%%%%%%%%%%%%%

\pagestyle{empty}
\begin{abstract}

Historicamente, a bioloxía mariña dependeu da inspección visual de expertos para avaliar a saúde
dos ecosistemas. No caso das algas laminarias (Laminaria hyperborea), a análise de depredación por
peixes e invertebrados realízase mediante a asignación manual de notas nunha escala numérica de
enteiros, un proceso subxectivo e custoso en tempo ante o crecente volume de datos visuais
dispoñibles. Concretamente neste caso, non existe estudos previos que traten
específicamente este tipo de alga. Ademáis, existe un problema de preprocesado e normalización
que solucionar nas imaxes, debido a que se tomaron de unha forma non estandarizada.Este proxecto
aborda a automatización de citado proceso mediante un pipeline de dúas etapas de aprendizaxe
profundo. Utilízase un dataset de 4.977 imaxes proporcionado pola Facultade de Ciencias para
adestrar, en primeira instancia, un modelo YOLO encargado da detección e localización das algas nas
imaxes. Posteriormente, implementase unha arquitectura de regresión ordinal (CORAL-CNN) para
clasificar o nivel de dano. A diferenza da clasificación tradicional, este enfoque preserva a orde lóxica
da etiquetación por dano incremental, permitindo que o sistema funcione como unha ferramenta
robusta de soporte á decisión para os expertos, reducindo a subxectividade e optimizando o fluxo de
traballo científico

  \vspace*{25pt}
  \begin{segundoresumo}
  
\textcolor{red}{\textbf{TODO: Simplemente traducir o resumo en galego e inglés.}}


  \end{segundoresumo}
\vspace*{25pt}
\begin{multicols}{2}
\begin{description}
\item [\palabraschaveprincipal:] \mbox{} \\[-20pt]
  \begin{itemize}
      \item YOLO
      \item CNN
      \item Visión por computador
      \item CORAL-CNN
      \item Alga Laminaria
      \item HPI
      \item IVR
  \end{itemize}
\end{description}
\begin{description}
\item [\palabraschavesecundaria:] \mbox{} \\[-20pt]
  \begin{itemize}
      \item YOLO
      \item CNN
      \item Computer Vision
      \item CORAL-CNN
      \item Kelp
      \item HPI
      \item IVR
  \end{itemize}
\end{description}
\end{multicols}

\end{abstract}
\pagestyle{fancy}

%%%%%%%%%%%%%%%%%%%%%%%%%%%%%%%%%%%%%%%%%%%%%%%%%%%%%%%%%%%%%%%%%%%%%%%%%%%%%%%%
